\section{Number of NGEs to the Right}
\label{sec:sq-num-nges}

\problemstatement{Given an array $\textit{arr}$ of $n$ elements and a list
of query indices, for each query index $i$, count \textbf{how many}
elements to the right of $i$ (not just the nearest one) are strictly
greater than $\textit{arr}[i]$.}

\begin{patternbox}
This problem is grouped alongside the monotonic-stack problems because it
shares their right-to-left scanning direction, but the keyword
\emph{``count how many,''} rather than \emph{``find the nearest,''} is a
signal that a plain monotonic stack is not quite the right tool: a stack
only ever keeps a thin, decreasing/increasing \emph{skyline} of values, not
the full multiset of everything seen so far, so it cannot directly answer
``how many are greater than $X$.'' This calls for a different right-to-left
technique -- a \textbf{Fenwick tree (Binary Indexed Tree) over compressed
ranks} -- that still processes the array in the same right-to-left order
as the rest of this section, but tracks counts rather than a skyline.
\end{patternbox}

\subsection*{Understanding the problem carefully}

Scanning right to left, by the time we reach index $i$, we want to know
how many of the values already seen (i.e.\ to $i$'s right) exceed
$\textit{arr}[i]$. This is a classic ``count of elements greater than $X$
among those inserted so far'' query, repeated $n$ times with the
collection growing by one element after each query -- exactly what a
Fenwick tree, indexed by the \emph{rank} of each value (after coordinate
compression, since values may be arbitrary integers), is built for: it
supports both ``insert a value'' and ``count how many inserted values are
$\le$ some threshold'' in $O(\log n)$ each.

\begin{ideabox}[Fenwick Tree over Compressed Ranks, Scanned Right to Left]
Compress all values in $\textit{arr}$ to ranks $1, \dots, n$ (sorted
order). Scan $\textit{arr}$ from right to left. For each $\textit{arr}[i]$:
query the Fenwick tree for the count of already-inserted values with rank
\emph{greater than} $\textit{arr}[i]$'s rank (equivalently, $n$ minus the
count of values with rank $\le \textit{arr}[i]$'s rank) -- this is exactly
the number of elements to $i$'s right that exceed $\textit{arr}[i]$.
Record this count for index $i$. Then insert $\textit{arr}[i]$'s rank into
the Fenwick tree, so it is correctly counted for any index further to the
left.
\end{ideabox}

\subsection*{Why scanning right to left and inserting afterward is correct}

At the moment we query for index $i$, the Fenwick tree contains exactly
the ranks of elements at indices $i+1, \dots, n-1$ -- because we always
insert $\textit{arr}[i]$ \emph{after} querying for it, and we process
indices in strictly decreasing order, so every index to $i$'s right has
already been inserted, and no index at or to $i$'s left has been inserted
yet. This invariant is what makes the Fenwick tree's count query, at the
moment we use it for index $i$, exactly answer ``how many elements to the
right of $i$ exceed $\textit{arr}[i]$,'' with no extra bookkeeping needed.

\subsection*{Algorithm}

\begin{enumerate}
  \item Compute ranks of all values in $\textit{arr}$ via sorting
        (coordinate compression).
  \item Initialize an empty Fenwick tree of size $n$.
  \item For $i$ from $n-1$ down to $0$:
  \begin{enumerate}
    \item $\textit{countGreater}[i] \gets n - \Call{Query}{\textit{rank}[i]}$
          (count of inserted values with rank $\le \textit{rank}[i]$,
          subtracted from the total inserted so far).
    \item $\Call{Update}{\textit{rank}[i], +1}$ (insert this rank).
  \end{enumerate}
  \item Answer each query index by direct lookup into
        $\textit{countGreater}$.
\end{enumerate}

\subsection*{Dry run}

\dryrun{Take $\textit{arr} = [3,4,2,7,5]$, query index $i=2$ (value
$2$).}

Ranks (1-indexed, smallest $=1$): $3\to2, 4\to3, 2\to1, 7\to5, 5\to4$.

Scanning right to left and inserting:

\begin{itemize}
  \item $i=4$ (value $5$, rank $4$): tree empty, query returns $0$
        inserted so far, so $\textit{countGreater}[4]=0-0=0$. Insert
        rank $4$.
  \item $i=3$ (value $7$, rank $5$): inserted so far $=1$ (rank $4$),
        count with rank $\le5$ is $1$. $\textit{countGreater}[3]=1-1=0$.
        Insert rank $5$.
  \item $i=2$ (value $2$, rank $1$): inserted so far $=2$ (ranks
        $4,5$), count with rank $\le1$ is $0$.
        $\textit{countGreater}[2]=2-0=2$. Insert rank $1$.
\end{itemize}

$\textit{countGreater}[2]=2$: indeed, to the right of index $2$ (value
$2$), the elements are $7$ and $5$, both greater -- $2$ elements,
matching.

\subsection*{C++ implementation}

\begin{lstlisting}
#include <bits/stdc++.h>
using namespace std;

class FenwickTree {
    vector<int> tree;
    int n;
public:
    FenwickTree(int n) : n(n), tree(n + 1, 0) {}

    void update(int i, int delta) {
        for (; i <= n; i += i & (-i)) tree[i] += delta;
    }

    int query(int i) { // count of inserted values with rank <= i
        int sum = 0;
        for (; i > 0; i -= i & (-i)) sum += tree[i];
        return sum;
    }
};

vector<int> countGreaterToRight(vector<int>& arr) {
    int n = arr.size();
    vector<int> sorted = arr;
    sort(sorted.begin(), sorted.end());

    auto rankOf = [&](int val) {
        return lower_bound(sorted.begin(), sorted.end(), val) - sorted.begin() + 1;
    };

    FenwickTree ft(n);
    vector<int> countGreater(n);
    int insertedSoFar = 0;

    for (int i = n - 1; i >= 0; i--) {
        int r = rankOf(arr[i]);
        countGreater[i] = insertedSoFar - ft.query(r);
        ft.update(r, 1);
        insertedSoFar++;
    }
    return countGreater;
}

int main() {
    vector<int> arr = {3, 4, 2, 7, 5};
    auto result = countGreaterToRight(arr);
    for (int x : result) cout << x << " ";
    cout << endl;
    return 0;
}
\end{lstlisting}

\begin{complexitybox}
\textbf{Time Complexity.} Coordinate compression costs $O(n \log n)$. The
main scan performs $n$ Fenwick tree updates and queries, each
$O(\log n)$, giving an overall time complexity of
\[
  O(n \log n).
\]
\textbf{Space Complexity.} $O(n)$ for the Fenwick tree and rank arrays.
\end{complexitybox}

\begin{takeawaybox}
``Count how many'' is a structurally different question from ``find the
nearest,'' even when both scan in the same direction -- a monotonic stack
answers the latter in $O(n)$ total but cannot answer the former at all,
since it discards the very elements a counting query would need. Whenever
a problem asks for a \emph{count} of qualifying elements rather than the
nearest one, reach for a Fenwick tree, segment tree, or balanced
order-statistics structure instead of a plain stack.
\end{takeawaybox}
